% Version 3.10: preserves v3.9 public import; adapts only its selection interface.
% Insert after the installed-snapshot relay and before INPUT_REPORT.
\paragraph{Request import from retained public DATA.}
\label{sec:public-request-import}
Each correct node's persistent request owner also runs the local import
loop below. Its input is the actual bytes already retained in
\tagg{PublicStore}, not a request query to an old roster. Its output is
a store of unvalidated request candidates and, after the ordinary
request checks, entries in the existing pending FIFO. Import neither
validates its containing DATA nor changes $G$, a boundary, a lock, a
vote, a report or an entry gate. In particular, it does not reopen a
closed session.

The owner retains an append-only source-work journal $Q^{\rm src}_v$,
an append-only candidate store $U_v$, and an append-only checking journal
$Q^{\rm chk}_v$. Each journal has a durable completion cursor and, for
its current finite job, a saved local subposition. Candidate entries
retain their exact bytes, source-object occurrence paths and local
arrival numbers. Checking outcomes and pending-store updates are
durable. A completed job is marked complete; its evidence is not
deleted. Existing pending entries retain their first-receipt order.

When complete raw source content is retained, append its actual retained
root and dictionary cut to $Q^{\rm src}_v$ before releasing the source
input or its enqueue obligation. Do the same when an exact public
reference target is subsequently acquired. On first enabling this rule,
enqueue one fixed finite cut of the already retained roots; concurrent
acquisitions append separate jobs. This is a local scan of held bytes,
not a new hydration or Ready condition. The new journal cursor is
independent of source availability, P6 and regional entry cursors.
Retaining the descriptor does not run its scan in a timed handler.

A source job follows the original typed fields and list positions of
\cref{sec:logical-carrier}, including DATA in proposals, ancestry and
boundary evidence. In the existing block tuple
\[
 B=(\beta,r,k,h,\mathit{type},\mathit{batch},
          \Hash(B_{\rm parent}),\rho,\xi),
\]
it extracts the actual entries of $\mathit{batch}$ only when the local
structure identifies a DATA body and supplies those entry bytes.
DRAIN has no batch. Arbitrary witness strings, request hashes and a
sender's object index are not DATA batch entries. Decode independently
represented typed children even if a sibling or enclosing proof fails.
Malformed bytes without that logical structure yield no invented entry.

The job traverses supplied inline objects and exact public targets held
at its fixed dictionary cut, visiting each actual object once. It saves
an unresolved edge without waiting for or fetching its target; a later
acquisition supplies another job. A cycle cannot cause repeated visits.
Thus the job has a finite fixed input. It does not fetch AUTH, native
state, source history or optional proof targets. Outer chain, version,
base and interval hints, P6 rejection and branch closure do not filter
this scan. An actual DATA entry is still only a candidate, including
when the containing DATA claims the wrong signed domain.

Append each new exact candidate byte string to $U_v$ before advancing
past its batch position. Before successful request validation,
deduplicate only identical bytes, with byte equality or a binding
content commitment, and preserve their source occurrences. A claimed
client identifier creates no duplicate marker. Invalid encodings and
failed authentication cannot reserve an identifier or suppress a later
valid request. $U_v$ is explicitly unvalidated: its entries are not
eligible reports, authenticated relays or executable requests.

\paragraph{Checking candidates without a moving target.}
When the ordinary application interface has completely prepared an exact
parent $P$, enqueue a checking job for the finite current cut of $U_v$.
When a new candidate is retained and such a parent is already prepared,
enqueue a check of that candidate at that parent. Equal
candidate-byte/parent checks need not be repeated. Each job retains the
already available read context needed by the ordinary finite request
check, including its exact native state, until the check finishes.
A pointer to a subsequently mutated engine state is not such a context.
Any required snapshot construction must actually finish before this
notification; neither its cost nor instantaneous copying is inferred
from the import rule. Ready does not wait for this construction.
Changing the active parent does not retarget or cancel an old job.
These jobs acquire no absent validation input and cause no state
reconstruction; preparing a parent keeps its original prerequisites.

At that fixed $P$, run the ordinary chain/domain, client authentication,
canonicalization and application-admission checks. A failed check is an
outcome only for those bytes at that parent, not a permanent rejection
of an asserted identity or a substitute for checks at later parents.
A success establishes $\operatorname{enc}(x)$ and the request's actual
logical identity under the original request interface. Only now merge
retransmissions by that validated identity; an asserted equal identifier
alone is insufficient. If its identity is already in $G$, record the
outcome without adding an active pending entry. Otherwise persist it in
the ordinary pending store, preserving an existing entry's FIFO position
or assigning the next first-receipt position to a new entry, and invoke
the same installed-snapshot relay rule. A successful check at an old
parent is not admission at a newer one.

The loop continuously alternates finite local steps on the oldest
unfinished source job and the oldest unfinished checking job, skipping
an empty journal. Within each job it follows the saved field or
candidate order. New jobs go to the tail and cannot overtake a fixed
job. The node continues these steps until each queued finite job is
complete; session changes and repeated new arrivals cannot suspend it
indefinitely. This is an explicit eventual local-work duty, with no
wall-clock deadline, throughput bound or finite-storage claim.
Neither queue waits on a fetch or on the other queue's completion.
Ordinary messages, timers, P6, Ready and report emission do not wait for
the import loop, its backlog, or an import watermark. Their original
processing bounds are not bounds on this deferred work.

The journals, bytes, read contexts and pending enqueue obligations belong
to the node, independently of a producing engine or import worker.
Waiting, Frozen, Completed, Terminated, \texttt{RequestClose}, Closed,
checkpoint compaction and advancement of $G$ cannot drop them. A worker
exit transfers them durably and a restart resumes their exact cursors.
The existing public-source retention rules continue to govern their
underlying bytes. Inclusion in $G$ removes an active pending request;
it does not erase raw evidence or turn a rejected record into accepted
history.

\paragraph{Interface to prepared input and FIFO filling.}
The import mechanism remains deferred and nonblocking. Its unvalidated
$U_v$ entries are not authenticated pending entries, offers or executable
requests. A successful check enters the persistent canonical/pending
store and the same request-receipt/relay hook. The selection job of
\cref{sec:input-selection} freezes a raw-journal/canonical-store cut after
hydration, then actually completes its own exact-opening admission and
minimum computation before signing an offer. Raw complete request
handles delivered by hydration are included even if this independent
public-import loop has not yet canonicalized a matching occurrence.
An unscanned public source root alone is not an invented request handle.

REPORT validates the supplied opening certificate and initializer choice,
or acknowledges a discharged ancestry obligation; it does not inspect
this import backlog or compute a live pending minimum. An import after
a frozen cut cannot amend that cut, its offer or an emitted report.
It remains available to later preparations and ordinary FIFO filling.
For FIFO, validate the actual parent, its native state and complete
ancestry, including inherited and speculative DATA, and recheck admission
after earlier batch entries. A previous import check or public occurrence
replaces none of these checks.

\begin{proposition}[Import after a first public carrier]
\label{prop:v39-public-request-import}
Suppose an original complete public result $R_0$ contains the complete
bytes of a request $x$ in an actual DATA batch in its old logical fields
or unchanged public source denotations. If a correct node follows and
retains those actual bytes and continues the local loop, it eventually
retains $x$ in $U_v$, even if the record is rejected or its branch is
closed. If, after that retention, an exact parent at which $x$ passes
the ordinary request checks is completely prepared, or such a parent
is already prepared then, the request eventually enters the pending
store unless its identity has already entered $G$.
\end{proposition}
\begin{proof}
The original public fields determine the DATA body and its actual batch
entries. Source retention records a job for each acquired root or exact
target. Each fixed job has only finitely many predecessors, a finite
input and finite local traversal; tail arrivals cannot overtake it.
The prescribed loop therefore reaches the actual entry and persists its
bytes. Failed hints, branch closure and unrelated witness failures are
not scan gates. An unresolved edge alone supplies no bytes, but the
hypothesis supplies acquisition of its actual target and hence its job.
Exact-byte deduplication only finds a copy already retained; an invalid
claimed identifier cannot suppress the entry.

The prepared-parent event or candidate arrival likewise records a fixed
finite checking job. Its retained local context needs no fetch. FIFO
service eventually completes the ordinary successful check of $x$ and
updates the pending store. A later parent cannot cancel that work.
Thereafter the ordinary pending rule retains the request until its
identity enters $G$, while each later selection job and batch independently checks
eligibility at its own exact parent.
\end{proof}

This result gives no bound from public visibility to pending admission.
In particular, valid candidate bytes cannot be omitted forever merely
by leaving their finite checking job unprocessed, but import need not
finish before any particular report. The proposition neither forms
the first $x$-bearing $R_0$ nor ensures that every future roster knows
$x$ before its first DATA. It adds no request-publication service,
cross-snapshot message guarantee or historical quorum-availability
premise. The ownership, finite preparation and session-first argument in
\cref{prop:session-input-progress} provide the separate route from an
actual correct-owner receipt to inclusion under continued regional service.
